% Auto-generated by experiments/19_report_visualization_enrichment.py
\section*{Visual Evidence Gallery}
\begin{figure}[!htbp]
\centering
\begin{subfigure}[t]{0.48\linewidth}
\centering
\includegraphics[width=\linewidth]{results/figures/report_enrichment/stage02_solver_validation_tridiagonal_operator.pdf}
\caption{The discretized Black Scholes operator has the expected sparse tridiagonal structure.}
\end{subfigure}
\hfill
\begin{subfigure}[t]{0.48\linewidth}
\centering
\includegraphics[width=\linewidth]{results/figures/report_enrichment/stage02_solver_validation_fd_stencil_validation.pdf}
\caption{The validation audit checks that the implementation uses the intended three point stencil.}
\end{subfigure}
\end{figure}

\begin{figure}[!htbp]
\centering
\begin{subfigure}[t]{0.48\linewidth}
\centering
\includegraphics[width=\linewidth]{results/figures/report_enrichment/stage02_solver_validation_coefficient_sign_audit.pdf}
\caption{The coefficient diagram helps connect operator sign conventions to implemented matrix rows.}
\end{subfigure}
\hfill
\begin{subfigure}[t]{0.48\linewidth}
\centering
\includegraphics[width=\linewidth]{results/figures/report_enrichment/stage02_solver_validation_annotated_lh_matrix.pdf}
\caption{The labeled operator matrix makes the lower diagonal and upper entries explicit for implementation review.}
\end{subfigure}
\end{figure}

\begin{figure}[!htbp]
\centering
\begin{subfigure}[t]{0.48\linewidth}
\centering
\includegraphics[width=\linewidth]{results/figures/report_enrichment/stage02_solver_validation_cn_a_b_matrices.pdf}
\caption{The Crank--Nicolson matrix pair makes the A and B sign convention explicit for the Stage 2 implementation audit.}
\end{subfigure}
\hfill
\begin{subfigure}[t]{0.48\linewidth}
\centering
\includegraphics[width=\linewidth]{results/figures/report_enrichment/stage02_solver_validation_rhs_boundary_terms.pdf}
\caption{The right hand side diagram shows why boundary conditions must be audited together with matrix construction.}
\end{subfigure}
\end{figure}

\begin{figure}[!htbp]
\centering
\begin{subfigure}[t]{0.48\linewidth}
\centering
\includegraphics[width=\linewidth]{results/figures/report_enrichment/stage02_solver_validation_psor_projection_step.pdf}
\caption{The projected update enforces the payoff obstacle after each relaxed candidate update.}
\end{subfigure}
\hfill
\begin{subfigure}[t]{0.48\linewidth}
\centering
\includegraphics[width=\linewidth]{results/figures/report_enrichment/stage02_solver_validation_lcp_matrix_system.pdf}
\caption{The LCP diagram connects the CN matrices to obstacle and complementarity diagnostics.}
\end{subfigure}
\end{figure}

\begin{figure}[!htbp]
\centering
\begin{subfigure}[t]{0.48\linewidth}
\centering
\includegraphics[width=\linewidth]{results/figures/report_enrichment/stage02_solver_validation_time_step_flow.pdf}
\caption{Time marching starts from the nonsmooth payoff and builds the American value surface step by step.}
\end{subfigure}
\hfill
\begin{subfigure}[t]{0.48\linewidth}
\centering
\includegraphics[width=\linewidth]{results/figures/report_enrichment/stage02_solver_validation_validation_ladder.pdf}
\caption{Validation proceeds from analytic checks to inequality and shape diagnostics before downstream experiments.}
\end{subfigure}
\end{figure}

\begin{figure}[!htbp]
\centering
\begin{subfigure}[t]{0.48\linewidth}
\centering
\includegraphics[width=\linewidth]{results/figures/report_enrichment/stage02_solver_validation_smoke_error_comparison.pdf}
\caption{Closed-form European put and theorem-based no-dividend call checks quantify smoke-grid error before downstream classical experiments.}
\end{subfigure}
\hfill
\begin{subfigure}[t]{0.48\linewidth}
\centering
\includegraphics[width=\linewidth]{results/figures/report_enrichment/stage02_solver_validation_obstacle_complementarity.pdf}
\caption{Complementarity residuals remain near solver tolerance across the smoke cases; obstacle violation is zero for every case and is stated explicitly rather than shown as an empty bar chart.}
\end{subfigure}
\end{figure}

\FloatBarrier
